\chapter{相对论力学基础}

\section{自由粒子的相对论作用量}

在经典力学中，质点的运动由轨道 $q^{i}(t)$ 描述，其作用量写为
\[
S=\int_{t_{1}}^{t_{2}}L(q^{i},\dot q^{i},t)\,\dif t .
\]
真实轨道满足 Hamilton 原理
\[
\delta S=0 .
\]
这里的变分不是把端点也一起移动，而是在端点固定的条件下比较所有可能轨道。也就是说，真实轨道附近的任意一条试探轨道都可写成
\[
q^{i}(t)+\delta q^{i}(t),
\]
其中 $\delta q^{i}(t)$ 是任意小函数，但在初末时刻必须为零。
对轨道作无穷小变分
\[
q^{i}(t)\longrightarrow q^{i}(t)+\delta q^{i}(t),
\]
并取固定端点条件
\[
\delta q^{i}(t_{1})=\delta q^{i}(t_{2})=0,
\]
可得到通常的 Euler--Lagrange 方程：
\[
\frac{\dif}{\dif t}\left(\frac{\partial L}{\partial \dot q^{i}}\right)
-
\frac{\partial L}{\partial q^{i}}
=0 .
\]
其推导的关键一步是
\[
\delta\dot q^{i}
=\frac{\dif}{\dif t}\delta q^{i}.
\]
因此
\[
\delta S
=
\int_{t_{1}}^{t_{2}}
\left(
\frac{\partial L}{\partial q^{i}}\delta q^{i}
+
\frac{\partial L}{\partial \dot q^{i}}\delta\dot q^{i}
\right)\dif t .
\]
对第二项分部积分得
\[
\delta S
=
\left.
\frac{\partial L}{\partial \dot q^{i}}\delta q^{i}
\right|_{t_{1}}^{t_{2}}
+
\int_{t_{1}}^{t_{2}}
\left[
\frac{\partial L}{\partial q^{i}}
-
\frac{\dif}{\dif t}
\left(\frac{\partial L}{\partial \dot q^{i}}\right)
\right]\delta q^{i}\,\dif t .
\]
由于端点变分为零，边界项消失；又由于 $\delta q^{i}$ 在中间时刻任意，括号中的系数必须为零。这就得到 Euler--Lagrange 方程。

在相对论中，作用量不仅要给出正确的运动方程，还必须是 Lorentz 标量。对于自由粒子而言，世界线上最自然的 Lorentz 标量是固有时积分，或者等价地，是世界线长度积分。类时世界线满足
\[
\dif s^{2}=\eta_{\mu\nu}\dif x^{\mu}\dif x^{\nu}=c^{2}\dif\tau^{2}.
\]
因此，自由粒子的相对论作用量取为
\[
S=-mc\int \dif s=-mc^{2}\int \dif\tau .
\]
这里 $m$ 为粒子的静质量。负号是一种约定，它保证低速极限下得到通常的经典自由粒子拉格朗日量。

若用任意参数 $\lambda$ 描述世界线，
\[
x^{\mu}=x^{\mu}(\lambda),
\]
则
\[
\dif x^{\mu}=\frac{\dif x^{\mu}}{\dif\lambda}\dif\lambda.
\]
于是作用量可写为
\[
S=-mc\int
\sqrt{\eta_{\mu\nu}
	\frac{\dif x^{\mu}}{\dif\lambda}
	\frac{\dif x^{\nu}}{\dif\lambda}}
\,\dif\lambda .
\]
如果换用另一个单调参数 $\lambda'=\lambda'(\lambda)$，则
\[
\frac{\dif x^{\mu}}{\dif\lambda}
=
\frac{\dif x^{\mu}}{\dif\lambda'}
\frac{\dif\lambda'}{\dif\lambda}.
\]
根号中的量会多出因子 $(\dif\lambda'/\dif\lambda)^{2}$，而积分测度 $\dif\lambda$ 同时变为
\[
\dif\lambda
=
\frac{\dif\lambda}{\dif\lambda'}\dif\lambda' .
\]
二者正好抵消。因此该形式不依赖参数 $\lambda$ 的具体选取，说明世界线作用量的重参数化不变性。

\section{自由粒子的相对论拉格朗日量}

在某一惯性系中，固有时与坐标时间满足
\[
\dif\tau=\dif t\sqrt{1-\frac{v^{2}}{c^{2}}},
\]
其中
\[
v^{2}=\bm{v}\cdot\bm{v}.
\]
这是由时空间隔
\[
c^{2}\dif\tau^{2}
=
c^{2}\dif t^{2}-\dif\bm{x}^{2}
\]
直接得到的。因为 $\dif\bm{x}=\bm{v}\dif t$，所以
\[
c^{2}\dif\tau^{2}
=
c^{2}\dif t^{2}-v^{2}\dif t^{2}
=
c^{2}\dif t^{2}
\left(1-\frac{v^{2}}{c^{2}}\right).
\]
即
\[
\dif\tau
=
\dif t
\sqrt{1-\frac{v^{2}}{c^{2}}}.
\]
将其代入自由粒子作用量，得到
\[
S=-mc^{2}\int \dif\tau
=-mc^{2}\int \sqrt{1-\frac{v^{2}}{c^{2}}}\,\dif t .
\]
与 $S=\int L\,\dif t$ 比较，自由粒子的相对论拉格朗日量为
\[
L=-mc^{2}\sqrt{1-\frac{v^{2}}{c^{2}}}.
\]

为了检查其非相对论极限，令 $v\ll c$。利用展开式
\[
\sqrt{1-\frac{v^{2}}{c^{2}}}
=
1-\frac{1}{2}\frac{v^{2}}{c^{2}}
+O\left(\frac{v^{4}}{c^{4}}\right),
\]
得到
\[
L=-mc^{2}+\frac{1}{2}mv^{2}
+O\left(\frac{v^{4}}{c^{2}}\right).
\]
第一项 $-mc^{2}$ 是常数，不影响 Euler--Lagrange 方程。因此低速极限下，上式回到经典自由粒子的拉格朗日量
\[
L_{\mathrm{nr}}=\frac{1}{2}mv^{2}.
\]

\section{自由粒子的世界线方程}

自由粒子的世界线方程可以直接由作用量变分得到。取
\[
\dot x^{\mu}=\frac{\dif x^{\mu}}{\dif\lambda},
\qquad
\dot x^{2}=\eta_{\mu\nu}\dot x^{\mu}\dot x^{\nu}.
\]
作用量写为
\[
S=-mc\int \sqrt{\dot x^{2}}\,\dif\lambda .
\]
相应的有效拉格朗日量为
\[
L_{\lambda}=-mc\sqrt{\eta_{\mu\nu}\dot x^{\mu}\dot x^{\nu}}.
\]
这里的 $L_{\lambda}$ 不是以坐标时间为自变量的普通拉格朗日量，而是以世界线参数 $\lambda$ 为自变量的拉格朗日量。变分时把 $x^{\mu}(\lambda)$ 看成四个广义坐标，端点条件为
\[
\delta x^{\mu}(\lambda_{1})=
\delta x^{\mu}(\lambda_{2})=0 .
\]
由于 $L_{\lambda}$ 不显含 $x^{\mu}$，Euler--Lagrange 方程给出
\[
\frac{\dif}{\dif\lambda}
\left(
\frac{\partial L_{\lambda}}{\partial \dot x^{\mu}}
\right)=0 .
\]
直接计算有
\[
\frac{\partial L_{\lambda}}{\partial \dot x^{\mu}}
=
-mc\frac{\dot x_{\mu}}{\sqrt{\dot x^{2}}}.
\]
其中用到了
\[
\frac{\partial\dot x^{2}}{\partial\dot x^{\mu}}
=
\frac{\partial}{\partial\dot x^{\mu}}
(\eta_{\alpha\beta}\dot x^{\alpha}\dot x^{\beta})
=2\eta_{\mu\nu}\dot x^{\nu}
=2\dot x_{\mu}.
\]
因此
\[
\frac{\dif}{\dif\lambda}
\left(
\frac{\dot x_{\mu}}{\sqrt{\dot x^{2}}}
\right)=0 .
\]
取固有时参数 $\lambda=\tau$ 时，
\[
\dot x^{\mu}=\frac{\dif x^{\mu}}{\dif\tau}=u^{\mu},
\qquad
\dot x^{2}=u^{\mu}u_{\mu}=c^{2}.
\]
于是上式化为
\[
\frac{\dif u_{\mu}}{\dif\tau}=0,
\]
等价地，
\[
\frac{\dif^{2}x^{\mu}}{\dif\tau^{2}}=0 .
\]
这说明自由粒子的世界线是闵可夫斯基时空中的直线。它是相对论中惯性运动的几何表述。

\section{四速度与四加速度}

粒子的四维坐标写为
\[
x^{\mu}=(ct,\bm{x}).
\]
固有时 $\tau$ 是类时世界线的自然参数。四速度定义为
\[
u^{\mu}=\frac{\dif x^{\mu}}{\dif\tau}.
\]
注意四速度不是三维速度 $\bm{v}$ 后面简单添一个分量；它是四维坐标对固有时的导数。由于固有时是 Lorentz 标量，$u^\mu$ 才按四维矢量变换。
由
\[
\dif t=\gamma\dif\tau,
\qquad
\gamma=\frac{1}{\sqrt{1-\frac{v^{2}}{c^{2}}}},
\]
可得时间分量
\[
u^{0}=\frac{\dif(ct)}{\dif\tau}=\gamma c,
\]
空间分量
\[
u^{i}=\frac{\dif x^{i}}{\dif\tau}=\frac{\dif x^{i}}{\dif t}\frac{\dif t}{\dif\tau}=\gamma v^{i}.
\]
因此
\[
u^{\mu}=\gamma(c,\bm{v}).
\]
四速度的 Minkowski 模长为
\[
u^{\mu}u_{\mu}=\gamma^{2}(c^{2}-v^{2})=c^{2}.
\]
这说明四速度的 Minkowski 长度恒为 $c$。

四加速度定义为
\[
a^{\mu}=\frac{\dif u^{\mu}}{\dif\tau}.
\]
虽然四速度的模长恒为 $c$，四速度本身仍可随固有时改变。四加速度描述的正是这种方向上的改变，而不是模长的改变。
对恒等式
\[
u^{\mu}u_{\mu}=c^{2}
\]
求固有时导数，得到
\[
\frac{\dif}{\dif\tau}(u^{\mu}u_{\mu})=2a^{\mu}u_{\mu}=0.
\]
因此
\[
a^{\mu}u_{\mu}=0 .
\]
四加速度总是与四速度正交。这里的正交是 Minkowski 意义下的正交。

\section{四动量与能量动量关系}

四动量定义为
\[
p^{\mu}=mu^{\mu}.
\]
这里假定粒子的质量 $m$ 为常数。因此四动量与四速度方向相同，并且也构成四维矢量。
由
\[
u^{\mu}=\gamma(c,\bm{v})
\]
得到
\[
p^{\mu}=(\gamma mc,\gamma m\bm{v}).
\]
定义相对论三维动量和能量为
\[
\bm{p}=\gamma m\bm{v},
\qquad
E=\gamma mc^{2}.
\]
这些定义不是任意记号。空间分量在低速极限 $\gamma\to 1$ 时回到 Newton 力学的 $\bm{p}=m\bm{v}$；时间分量乘以 $c$ 得到的量在后面会等于哈密顿量，因此自然解释为粒子总能量。
则四动量可写成
\[
p^{\mu}=\left(\frac{E}{c},\bm{p}\right).
\]
降低指标得到
\[
p_{\mu}=\left(\frac{E}{c},-\bm{p}\right).
\]

四动量的 Minkowski 模长为
\[
p^{\mu}p_{\mu}=m^{2}u^{\mu}u_{\mu}=m^{2}c^{2}.
\]
另一方面，
\[
p^{\mu}p_{\mu}=\left(\frac{E}{c}\right)^{2}-p^{2}.
\]
因此
\[
\frac{E^{2}}{c^{2}}-p^{2}=m^{2}c^{2}.
\]
两边乘以 $c^{2}$，得到能量动量关系
\[
E^{2}=p^{2}c^{2}+m^{2}c^{4}.
\]
对于无质量粒子，令 $m=0$，得到
\[
E=pc.
\]

\section{相对论哈密顿量}

自由粒子的三维正则动量由拉格朗日量
\[
L=-mc^{2}\sqrt{1-\frac{v^{2}}{c^{2}}}
\]
给出：
\[
\bm{p}=\frac{\partial L}{\partial\bm{v}}=\gamma m\bm{v}.
\]
具体地，
\[
\frac{\partial L}{\partial v_{i}}
=
-mc^{2}\frac{1}{2}
\left(1-\frac{v^{2}}{c^{2}}\right)^{-1/2}
\left(-\frac{2v_{i}}{c^{2}}\right)
=
\gamma mv_{i}.
\]
哈密顿量定义为
\[
H=\bm{p}\cdot\bm{v}-L.
\]
代入 $\bm{p}=\gamma m\bm{v}$ 与 $L=-mc^{2}/\gamma$，得到
\[
H=\gamma mv^{2}+\frac{mc^{2}}{\gamma}.
\]
利用
\[
\frac{1}{\gamma^{2}}=1-\frac{v^{2}}{c^{2}},
\]
可将第二项写成
\[
\frac{mc^{2}}{\gamma}
=
\gamma mc^{2}\frac{1}{\gamma^{2}}
=
\gamma mc^{2}\left(1-\frac{v^{2}}{c^{2}}\right)
=
\gamma mc^{2}-\gamma mv^{2}.
\]
可化简为
\[
H=\gamma mc^{2}=E.
\]
因此自由相对论粒子的哈密顿量等于总能量。

在哈密顿形式中，哈密顿量应写为动量的函数。由能量动量关系可得
\[
H(\bm{p})=\sqrt{p^{2}c^{2}+m^{2}c^{4}}.
\]
也可以直接从 $\bm{p}=\gamma m\bm{v}$ 消去速度。由
\[
p^{2}
=
\gamma^{2}m^{2}v^{2}
=
m^{2}c^{2}(\gamma^{2}-1),
\]
得到
\[
\gamma mc^{2}
=
\sqrt{p^{2}c^{2}+m^{2}c^{4}}.
\]
当 $p\ll mc$ 时，
\[
H=mc^{2}\sqrt{1+\frac{p^{2}}{m^{2}c^{2}}}
=
mc^{2}+\frac{p^{2}}{2m}-\frac{p^{4}}{8m^{3}c^{2}}+\cdots .
\]
第一项为静止能，第二项为经典动能，后续项为相对论修正。

\section{四动量的洛伦兹变换}

四动量是四维矢量，因此在 Lorentz 变换下满足
\[
p'^{\mu}=\Lambda^{\mu}{}_{\nu}p^{\nu}.
\]
因为 $p^{0}=E/c$，四动量的时间分量变换后再乘以 $c$ 就得到能量变换式；空间分量则直接给出三维动量分量的变换式。
考虑 $K'$ 相对于 $K$ 沿 $x$ 轴正方向以速度 $V$ 运动。若采用从 $K$ 到 $K'$ 的 Lorentz boost，则
\[
E'=\gamma_{V}(E-Vp_{x}),
\]
\[
p'_{x}=\gamma_{V}\left(p_{x}-\frac{V}{c^{2}}E\right),
\]
以及
\[
p'_{y}=p_{y},
\qquad
p'_{z}=p_{z}.
\]
其中
\[
\gamma_{V}=\frac{1}{\sqrt{1-\frac{V^{2}}{c^{2}}}}.
\]
反过来，从 $K'$ 到 $K$ 的变换为
\[
E=\gamma_{V}(E'+Vp'_{x}),
\]
\[
p_{x}=\gamma_{V}\left(p'_{x}+\frac{V}{c^{2}}E'\right),
\]
\[
p_{y}=p'_{y},
\qquad
p_{z}=p'_{z}.
\]
这些公式说明，能量和动量在 Lorentz 变换下相互混合。能量不是 Lorentz 标量，三维动量也不是孤立的四维对象；它们共同组成四动量。

\section{四维力与相对论运动方程}

三维力定义为
\[
\bm{f}=\frac{\dif\bm{p}}{\dif t}.
\]
相对论中更自然的对象是四维力，定义为
\[
F^{\mu}=\frac{\dif p^{\mu}}{\dif\tau}.
\]
由于
\[
\frac{\dif}{\dif\tau}=\gamma\frac{\dif}{\dif t},
\]
四维力的空间分量为
\[
F^{i}=\frac{\dif p^{i}}{\dif\tau}=\gamma\frac{\dif p^{i}}{\dif t}=\gamma f^{i}.
\]
时间分量为
\[
F^{0}=\frac{\dif}{\dif\tau}\left(\frac{E}{c}\right)
=\frac{\gamma}{c}\frac{\dif E}{\dif t}.
\]
功率公式可由能量动量关系推出。对
\[
E^{2}=p^{2}c^{2}+m^{2}c^{4}
\]
关于坐标时间求导，得
\[
2E\frac{\dif E}{\dif t}
=
2c^{2}\bm{p}\cdot\frac{\dif\bm{p}}{\dif t}.
\]
利用 $E=\gamma mc^{2}$ 与 $\bm{p}=\gamma m\bm{v}$，有
\[
\frac{c^{2}\bm{p}}{E}=\bm{v}.
\]
因此
\[
\frac{\dif E}{\dif t}
=
\bm{v}\cdot\frac{\dif\bm{p}}{\dif t}.
\]
由功率公式
\[
\frac{\dif E}{\dif t}=\bm{f}\cdot\bm{v},
\]
可得
\[
F^{0}=\frac{\gamma}{c}\bm{f}\cdot\bm{v}.
\]
因此
\[
F^{\mu}=\left(\frac{\gamma}{c}\bm{f}\cdot\bm{v},\gamma\bm{f}\right).
\]
四维运动方程为
\[
\frac{\dif p^{\mu}}{\dif\tau}=F^{\mu}.
\]
由于
\[
p^{\mu}p_{\mu}=m^{2}c^{2}
\]
为常数，对固有时求导得
\[
2p_{\mu}F^{\mu}=0.
\]
又因 $p^{\mu}=mu^{\mu}$，所以
\[
F^{\mu}u_{\mu}=0.
\]
四维力总是与四速度正交。这是相对论动力学对力的基本约束。

\section{四维角动量张量}

三维角动量为
\[
\bm{L}=\bm{r}\times\bm{p}.
\]
它的分量可以写成
\[
L^{i}=\epsilon^{ijk}x^{j}p^{k}.
\]
这个写法提示我们，角动量本质上来自位置和动量的反对称组合。到了四维时空中，时间分量也应与空间分量一起参与这种反对称组合。
在四维时空中，与角动量对应的对象是二阶反对称张量：
\[
M^{\mu\nu}=x^{\mu}p^{\nu}-x^{\nu}p^{\mu}.
\]
显然
\[
M^{\mu\nu}=-M^{\nu\mu}.
\]
其纯空间分量
\[
M^{ij}=x^{i}p^{j}-x^{j}p^{i}
\]
对应三维角动量。例如
\[
M^{12}=x^{1}p^{2}-x^{2}p^{1}=L^{3}.
\]
含时间指标的分量
\[
M^{0i}=x^{0}p^{i}-x^{i}p^{0}
\]
与 boost 生成元和质心运动相关。

对固有时求导：
\[
\frac{\dif M^{\mu\nu}}{\dif\tau}
=
\frac{\dif x^{\mu}}{\dif\tau}p^{\nu}
+x^{\mu}\frac{\dif p^{\nu}}{\dif\tau}
-
\frac{\dif x^{\nu}}{\dif\tau}p^{\mu}
-
x^{\nu}\frac{\dif p^{\mu}}{\dif\tau}.
\]
利用 $p^{\mu}=mu^{\mu}$，有
\[
u^{\mu}p^{\nu}-u^{\nu}p^{\mu}=m(u^{\mu}u^{\nu}-u^{\nu}u^{\mu})=0.
\]
因此
\[
\frac{\dif M^{\mu\nu}}{\dif\tau}
=
x^{\mu}F^{\nu}-x^{\nu}F^{\mu}.
\]
定义四维力矩张量
\[
N^{\mu\nu}=x^{\mu}F^{\nu}-x^{\nu}F^{\mu},
\]
则
\[
\frac{\dif M^{\mu\nu}}{\dif\tau}=N^{\mu\nu}.
\]
对于自由粒子，$F^{\mu}=0$，所以
\[
\frac{\dif M^{\mu\nu}}{\dif\tau}=0.
\]
这就是自由粒子的四维角动量张量守恒。

\section{哈密顿--雅可比方程}

Hamilton--Jacobi 理论把动力学问题转化为作用量函数的偏微分方程。在相对论中，可把作用量看成时空坐标的标量函数：
\[
S=S(x^{\mu}).
\]
这里的 $S(x^\mu)$ 可以理解为从某个固定初始点到事件 $x^\mu$ 的经典作用量。改变终点位置时，作用量的变化由终点处的动量决定；这正是 Hamilton--Jacobi 理论中把动量写成作用量梯度的原因。
为了与平面波相位
\[
S=-Et+\bm{p}\cdot\bm{x}
\]
相容，这里取约定
\[
p_{\mu}=-\partial_{\mu}S.
\]
这个负号来自本讲义对作用量相位的约定。若把 $S=-Et+\bm{p}\cdot\bm{x}$ 看成一条自由粒子轨道对应的作用量函数，则时间导数应给出能量，空间梯度应给出三维动量。
需要注意，$p_{\mu}$ 是四维协变动量分量，而不是三维动量分量。由于
\[
\partial_{\mu}=\left(\frac{1}{c}\frac{\partial}{\partial t},\bm{\nabla}\right),
\]
时间分量给出
\[
p_{0}=-\frac{\partial S}{\partial x^{0}}
=-\frac{1}{c}\frac{\partial S}{\partial t}.
\]
对于 $S=-Et+\bm{p}\cdot\bm{x}$，有
\[
p_{0}=\frac{E}{c}.
\]
空间分量为
\[
p_{i}=-\frac{\partial S}{\partial x^{i}}.
\]
在号差约定 $\eta_{\mu\nu}=\mathrm{diag}(1,-1,-1,-1)$ 下，四维协变空间分量满足 $p_{i}=-p^{i}$。因此三维动量为
\[
\bm{p}=\bm{\nabla}S.
\]

自由粒子的质量壳条件为
\[
p_{\mu}p^{\mu}=m^{2}c^{2}.
\]
代入 $p_{\mu}=-\partial_{\mu}S$，得到相对论 Hamilton--Jacobi 方程：
\[
\eta^{\mu\nu}\partial_{\mu}S\partial_{\nu}S=m^{2}c^{2}.
\]
展开为
\[
\frac{1}{c^{2}}\left(\frac{\partial S}{\partial t}\right)^{2}
-
(\bm{\nabla}S)^{2}
=
m^{2}c^{2}.
\]
若取 $S=-Et+\bm{p}\cdot\bm{x}$，则立即得到
\[
E^{2}=p^{2}c^{2}+m^{2}c^{4}.
\]
因此 Hamilton--Jacobi 方程可以看成能量动量关系在作用量函数上的局部形式。

\section{从粒子力学到场论}

本章讨论的是点粒子的相对论力学。点粒子的自由度是有限个世界线坐标
\[
x^{\mu}=x^{\mu}(\tau).
\]
场论的基本对象则是定义在时空每一点上的场，例如
\[
\phi(x),
\qquad
A^{\mu}(x),
\qquad
T^{\mu\nu}(x).
\]
因此，场论可以看成从有限自由度系统到无限自由度系统的推广。

粒子力学中的作用量通常是世界线上的积分：
\[
S=\int L_{\tau}\,\dif\tau.
\]
场论中的作用量则是时空区域上的积分：
\[
S=\int \mathcal{L}(\phi,\partial_{\mu}\phi,x)\,\dif^{4}x.
\]
其中 $\mathcal{L}$ 称为拉格朗日密度。

粒子力学中的 Euler--Lagrange 方程为
\[
\frac{\dif}{\dif t}
\left(
\frac{\partial L}{\partial \dot q^{i}}
\right)
-
\frac{\partial L}{\partial q^{i}}
=0.
\]
场论中的 Euler--Lagrange 方程将推广为
\[
\partial_{\mu}
\left(
\frac{\partial\mathcal{L}}{\partial(\partial_{\mu}\phi)}
\right)
-
\frac{\partial\mathcal{L}}{\partial\phi}
=0.
\]
由此可见，相对论力学不仅给出粒子的能量动量关系，也为经典场论提供了三个基本思想：作用量应是 Lorentz 标量；动力学方程应写成协变形式；连续对称性通过守恒量或守恒流体现。

\section{本章小结}

本章建立了相对论点粒子的基本力学形式。自由粒子的作用量为
\[
S=-mc\int \dif s=-mc^{2}\int \dif\tau,
\]
它是 Lorentz 标量。由该作用量可得到自由粒子在闵可夫斯基时空中的直线世界线。

四速度和四动量分别定义为
\[
u^{\mu}=\frac{\dif x^{\mu}}{\dif\tau},
\qquad
p^{\mu}=mu^{\mu}.
\]
它们满足
\[
u^{\mu}u_{\mu}=c^{2},
\qquad
p^{\mu}p_{\mu}=m^{2}c^{2}.
\]
由四动量的模长得到能量动量关系
\[
E^{2}=p^{2}c^{2}+m^{2}c^{4}.
\]

自由粒子的哈密顿量为
\[
H(\bm{p})=\sqrt{p^{2}c^{2}+m^{2}c^{4}},
\]
其低动量极限给出静止能、经典动能和相对论修正项。

四维力定义为
\[
F^{\mu}=\frac{\dif p^{\mu}}{\dif\tau},
\]
并满足约束
\[
F^{\mu}u_{\mu}=0.
\]
四维角动量张量为
\[
M^{\mu\nu}=x^{\mu}p^{\nu}-x^{\nu}p^{\mu}.
\]
自由粒子满足
\[
\frac{\dif p^{\mu}}{\dif\tau}=0,
\qquad
\frac{\dif M^{\mu\nu}}{\dif\tau}=0.
\]
这分别表示四动量守恒和四维角动量张量守恒。

最后，相对论 Hamilton--Jacobi 方程
\[
\eta^{\mu\nu}\partial_{\mu}S\partial_{\nu}S=m^{2}c^{2}
\]
表明能量动量关系也可以写成作用量函数的局部偏微分方程。

下一章将把这些思想推广到场论。粒子的世界线变量 $x^{\mu}(\tau)$ 将被时空场 $\phi(x)$、$A^{\mu}(x)$ 等取代；粒子力学中的守恒量将推广为场论中的守恒流和能量动量张量。
